\section{Propagation and the joint meromorphic boundary}
\label{sec:joint}

\subsection{A classical subharmonic propagation argument}

We give the precise local statement needed to pass from
Proposition~\ref{prop:slice} to a joint boundary. Its ingredients are
the standard upper-envelope and decreasing-limit principles for
subharmonic functions, Cauchy's inequality and the maximum principle;
see \citep[Section 4.8 and Properties 8.4.3]{kw2017}.
The lemma is an explicit application of these classical tools.

\begin{lemma}[Joint-cap propagation]\label{lem:cap}
Let \(V\subset\C\) be a connected domain and
\(F\in\mathcal O(\D\times V)\). Suppose that for one \(t\in V\),
the slice \(F(\,\cdot\,,t)\) has the whole unit circle as a
holomorphic natural boundary. Then \(F\) has no local joint
holomorphic or meromorphic cap through any point of \(\T\times V\).
\end{lemma}

\begin{proof}
Suppose first that there is a holomorphic cap \(H\) on
\(B(x_0,\rho)\times B(y_0,\eta)\), where \(|x_0|=1\) and
\(\overline{B(y_0,\eta)}\subset V\). It agrees with \(F\) on the
interior overlap; shrinking the product if necessary, the identity
theorem gives agreement throughout that connected overlap.
Choose
\[
0<\delta<\min(1/2,\rho/4),\qquad
z_c=(1-\delta)x_0,\qquad M=-\log\delta.
\]
For \(n\geq0\), the Taylor coefficients
\[
a_n(y)=\frac1{n!}\partial_x^nF(z_c,y)
\]
are holomorphic on \(V\).
For \(n\geq1\), let \(u_n(y)=n^{-1}\log|a_n(y)|\), with
\(\log0=-\infty\). Each \(u_n\) is subharmonic, allowing the
identically negative-infinite function.

For any \(0<q<\delta\) and compact \(K\subset V\), Cauchy's
inequality on the circle \(|x-z_c|=q\), uniformly for \(y\in K\),
gives a constant \(C_{q,K}\geq1\) such that
\begin{equation}\label{eq:cauchy-envelope}
u_n(y)\leq-\log q+\frac{\log C_{q,K}}n
\quad(y\in K,\ n\geq1).
\end{equation}
Thus the family is locally uniformly bounded above.
For a locally upper-bounded function \(h\), denote its upper
semicontinuous regularization by
\[
h^*(y)=\lim_{\varepsilon\downarrow0}
\sup_{|w-y|<\varepsilon}h(w).
\]
Define the regularized tail envelopes
\begin{equation}\label{eq:tail-envelopes}
V_N=\left(\sup_{n\geq N}u_n\right)^*,
\qquad V_\infty=\lim_{N\to\infty}V_N.
\end{equation}
The upper-envelope theorem makes each \(V_N\) subharmonic.
They decrease with \(N\), so \(V_\infty\) is subharmonic or
identically \(-\infty\).

Apply \eqref{eq:cauchy-envelope} on a compact neighborhood of an
arbitrary parameter point before regularizing. On its interior this
gives
\[
V_N\leq-\log q+\frac{\log C_{q,K}}N.
\]
Letting \(N\to\infty\) and then \(q\uparrow\delta\), we obtain
\begin{equation}\label{eq:global-envelope}
V_\infty(y)\leq M\qquad(y\in V).
\end{equation}
Regularization only raises the supremum; consequently
\[
V_\infty(t)\geq\limsup_{n\to\infty}u_n(t).
\]
The Taylor series of \(F(\,\cdot\,,t)\) about \(z_c\) has radius
exactly \(\delta\). It has at least that radius since
\(B(z_c,\delta)\subset\D\); a larger radius would continue the
slice through \(x_0\), contrary to its natural boundary.
Cauchy--Hadamard therefore gives
\(\limsup_nu_n(t)=-\log\delta=M\).
Together with \eqref{eq:global-envelope}, this yields
\(V_\infty(t)=M\). The strong maximum principle on connected \(V\)
forces \(V_\infty\equiv M\).

But \(\overline{B(z_c,2\delta)}\subset B(x_0,\rho)\), because
\(|x-x_0|\leq2\delta+\delta<\rho\) there.
For a smaller parameter disc \(B(y_0,\eta')\), the cap and Cauchy's
inequality now give
\[
|a_n(y)|\leq C'(2\delta)^{-n}
\quad(y\in\overline{B(y_0,\eta')}).
\]
These are the same coefficients, since \(H=F\) near \(z_c\).
Applying the tail-envelope argument on the interior of this disc
yields \(V_\infty\leq-\log(2\delta)<M\), a contradiction.
This excludes holomorphic caps.

For a meromorphic cap, shrink to a product where it can be written
as \(A/B\), with \(A,B\) holomorphic and \(B\not\equiv0\).
The denominator cannot vanish on an entire product of a unit-circle
arc and an open parameter disc. If it did, for every fixed parameter
the one-variable identity theorem in \(x\) would make \(B\)
identically zero on the first disc, and thus zero on the whole product.
There is therefore a nearby face point at which \(B\ne0\).
In a smaller neighborhood it supplies a holomorphic cap, already
excluded. This proves the meromorphic assertion.
\end{proof}

The regularized tails in \eqref{eq:tail-envelopes} avoid any assertion
that a raw pointwise logarithmic limsup is subharmonic, or that it
equals its regularization at every parameter.
The distinction is essential for exceptional slices: if
\(h\in\mathcal O(\D)\) has a natural boundary, then
\(F(x,y)=(y-y_0)h(x)\) has an entire zero slice at \(y_0\), yet still
has no joint cap anywhere on \(\T\times V\).

\subsection{The independent branch}

For multiplicatively independent bases,
Section~\ref{sec:convergence} proves holomorphy of \(R_{a,b}\) on
\(\D^2\). Choose \(t\in(0,1/b)\), apply
Proposition~\ref{prop:slice}, and then apply Lemma~\ref{lem:cap} with
\(V=\D\). This excludes every joint meromorphic cap on
\(\T\times\D\).
Since \(R_{a,b}(x,y)=R_{b,a}(y,x)\), the same argument excludes
\(\D\times\T\).
A cap at a corner in \(\T^2\) would contain a nearby point of one
of these faces and restrict to a cap there. Thus every point of
\(\partial(\D^2)\) is excluded in this branch.

\subsection{The dependent branch and its parameter threshold}

The dependent continuation has genuine poles in the bidisc, so one
must not apply Lemma~\ref{lem:cap} to it on all of \(\D^2\).
The coefficient bound \(C_{a,b}(n,m)\leq b^m-1\) does, however,
give joint holomorphy on
\[
\D\times\D_{1/b},\qquad \D_{1/b}=\{|y|<1/b\}.
\]
Proposition~\ref{prop:slice} and Lemma~\ref{lem:cap}, now with
\(V=\D_{1/b}\), exclude every joint meromorphic cap at a point
\((x_0,y_0)\) with \(|x_0|=1\) and \(|y_0|<1/b\).

For \(1/b<|y_0|<1\), we instead use the actual poles already
identified in Section~\ref{sec:poles}. Take arbitrarily large
integers \(k\) with \(\gcd(k,r)=1\). The pair \((sk,r)\) is primitive,
and
\[
h_{sk,r}=\gcd(rsk,sr)=rs.
\]
The slice poles include all the roots of
\begin{equation}\label{eq:boundary-pole-grids}
c^{rs}x^{sk}y_0^r=1.
\end{equation}
Their modulus is
\[
(c^{rs}|y_0|^r)^{-1/(sk)}<1,
\]
tending to one, and their complete root grids have mesh tending to
zero. Formula \eqref{eq:residue} proves that every such pole is
genuine, even if other components meet there.
Near each \(x_0\in\T\), infinitely many interior slice poles
therefore tend to \(x_0\).

Suppose a joint meromorphic cap \(A/B\) existed at such a
\((x_0,y_0)\). Shrink its parameter disc so that
\(1/b<|y|<1\) throughout it. For \emph{each} parameter \(y\) in
that disc, \eqref{eq:boundary-pole-grids} gives infinitely many
genuine slice poles in the cap's \(x\)-disc, tending to \(x_0\).
The denominator must vanish at each: otherwise \(A/B\) would be
holomorphic there and could not agree with the continued series.
The one-variable identity theorem makes \(B(\,\cdot\,,y)\)
identically zero on the \(x\)-disc. This for every \(y\) in an open
disc contradicts \(B\not\equiv0\).

If \(|y_0|=1/b\), any putative cap contains nearby face points with
\(|y|>1/b\), just excluded. The strip and pole arguments together
therefore exclude all of \(\T\times\D\).
Symmetry excludes the other face, and the same corner argument as
above excludes \(\T^2\).
This completes the proof of Theorem~\ref{thm:classification}.
